
\begin{abstract}

CodeRunner, the Moodle question type used to set and automatically grade programming exercises at the University of Strathclyde and elsewhere, is a mature and capable tool whose remaining weaknesses are problems of experience rather than of capability. Authoring a question means working through a dense editing form with fixed test-case batches, help text that vanishes on a click and misaligned controls, while its documentation is a single feature-organised page; students answer through a fixed layout with hidden feedback and little room to work. This project treated those frictions as an engineering problem and set out to remove them without disturbing what already worked.

The work proceeded from evidence to delivery. CodeRunner and its two audiences were analysed through heuristic review, interviews and a survey of students and teachers; the faults found were turned into prioritised requirements; and the improvements were built inside a fully reproducible Docker development environment that recreates the whole system with a single command. The delivered changes span both sides of the tool: flexible test-case management and a corrected Precheck display for authors, a switchable and resizable layout and reclaimable space for students, measurable accessibility gains, and task-focused documentation, with in-plugin search, served from within the plugin and reachable from the editing form.

The changes were verified by twenty-seven unit tests and the plugin's own browser-driven regression suite, all of which pass, and their value was confirmed by the survey: every delivered feature won majority approval from both students and teachers, no feature was rejected, and the accessibility audit rose from 93\% to 97\%. The significance of the result is methodological as much as practical: it demonstrates that a widely deployed educational tool can be improved meaningfully by disciplined user-experience engineering on the existing system, rather than by building a replacement.

\end{abstract}



